% ============================================================================
% LANGENTWURF: Ohmsches Gesetz - Einführung
% Physik Klasse 8, Gesamtschule heterogen (Profil A)
% ============================================================================

\documentclass[11pt,a4paper]{article}

% ============================================================================
% PACKAGES
% ============================================================================
\usepackage[utf8]{inputenc}
\usepackage[T1]{fontenc}
\usepackage[ngerman]{babel}
\usepackage{geometry}
\usepackage{fancyhdr}
\usepackage{lastpage}
\usepackage{xcolor}
\usepackage{tcolorbox}
\usepackage{enumitem}
\usepackage{tabularx}
\usepackage{longtable}
\usepackage{array}
\usepackage{booktabs}
\usepackage{hyperref}
\usepackage{ifthen}
\usepackage{tikz}
\usepackage{amssymb}
\usepackage{siunitx}

% SI-Einheiten
\sisetup{locale = DE}

% ============================================================================
% KONFIGURATION
% ============================================================================

% --- TIER (1 = Vollständig) ---
\newcommand{\plantier}{1}

% --- LERNGRUPPEN-PROFIL (A = Gesamtschule heterogen) ---
\newcommand{\lerngruppenprofil}{A}

% --- FACH ---
\newcommand{\fachid}{physik}

% Bedingte Befehle
\newcommand{\tierone}[1]{\ifthenelse{\equal{\plantier}{1}}{#1}{}}
\newcommand{\tieronetwo}[1]{\ifthenelse{\equal{\plantier}{1}\OR\equal{\plantier}{2}}{#1}{}}
\newcommand{\tierall}[1]{#1}

\newcommand{\testdifferenziert}[1]{%
    \ifthenelse{\equal{\fachid}{physik}\OR\equal{\fachid}{chemie}}{#1}{}%
}
\newcommand{\testeinheitlich}[1]{%
    \ifthenelse{\equal{\fachid}{informatik}\OR\equal{\fachid}{spanisch}}{#1}{}%
}

\newcommand{\profilheterogen}[1]{%
    \ifthenelse{\equal{\lerngruppenprofil}{A}}{#1}{}%
}
\newcommand{\profilhomogen}[1]{%
    \ifthenelse{\equal{\lerngruppenprofil}{D}\OR\equal{\lerngruppenprofil}{E}}{#1}{}%
}

% ============================================================================
% VARIABLEN
% ============================================================================

\newcommand{\fach}{Physik}
\newcommand{\jahrgangsstufe}{8}
\newcommand{\schulart}{Gesamtschule}
\newcommand{\kursart}{Regelunterricht}
\newcommand{\stundenthema}{Das Ohmsche Gesetz}
\newcommand{\reihenthema}{Elektrischer Widerstand}
\newcommand{\stundennummer}{1}
\newcommand{\reihenstunden}{4}
\newcommand{\unterrichtsdatum}{XX.XX.XXXX}
\newcommand{\unterrichtszeit}{XX:XX -- XX:XX}
\newcommand{\raum}{Physikraum}

\newcommand{\lehrkraft}{N.N.}
\newcommand{\ausbildungsstand}{LAA}

\newcommand{\lerngruppengroesse}{24}
\newcommand{\maennlich}{13}
\newcommand{\weiblich}{11}
\newcommand{\divers}{0}

\newcommand{\gerniveau}{}

\newcommand{\fachfarbe}{2c5aa0}

% ============================================================================
% LERNGRUPPEN-PROFIL
% ============================================================================

\newcommand{\profilbeschreibung}{%
    \textbf{Profil A: Gesamtschule heterogen}\\[0.5em]
    \begin{tabular}{ll}
        Klassenstufe: & 8 \\
        Zusammensetzung: & BR (20\%), MR (50\%), GY (30\%) \\
        Inklusion: & 2 SuS mit LRS, 1 SuS mit Förderbedarf Lernen \\
        Testdifferenzierung: & 3 Niveaus (*/**/***) \\
        Besonderheiten: & Hohe Heterogenität, intensive Differenzierung nötig
    \end{tabular}
}

% ============================================================================
% LAYOUT
% ============================================================================
\geometry{left=2.5cm, right=2.5cm, top=2.5cm, bottom=2.5cm}
\definecolor{fach}{HTML}{\fachfarbe}
\definecolor{gray}{HTML}{666666}
\definecolor{lightgray}{HTML}{f5f5f5}
\definecolor{successgreen}{HTML}{2a7a4b}
\definecolor{warningorange}{HTML}{b35c00}

\tcbuselibrary{skins,breakable}

\newtcolorbox{infobox}[1][]{
    colback=lightgray,
    colframe=fach,
    fonttitle=\bfseries,
    title=#1,
    breakable
}

\newtcolorbox{scorebox}{
    colback=white,
    colframe=fach,
    fonttitle=\bfseries,
    title=Didaktik-Score,
    breakable
}

\newtcolorbox{profilbox}{
    colback=lightgray,
    colframe=gray,
    fonttitle=\bfseries,
    title=Lerngruppen-Profil,
    breakable
}

\pagestyle{fancy}
\fancyhf{}
\fancyhead[L]{\small\textcolor{gray}{\fach{} -- Jg. \jahrgangsstufe{} -- \stundenthema}}
\fancyhead[R]{\small\textcolor{gray}{Langentwurf Tier \plantier{} / Profil \lerngruppenprofil}}
\fancyfoot[C]{\small\thepage{} / \pageref{LastPage}}
\renewcommand{\headrulewidth}{0.4pt}
\renewcommand{\footrulewidth}{0pt}

\renewcommand{\thesection}{\arabic{section}}
\renewcommand{\thesubsection}{\thesection.\arabic{subsection}}

% ============================================================================
% DOKUMENT
% ============================================================================
\begin{document}

% ============================================================================
% DECKBLATT
% ============================================================================
\begin{titlepage}
    \centering
    \vspace*{2cm}

    {\Large\textcolor{gray}{Greenhouse Schools Rostock}}\\[0.5cm]
    {\large\textcolor{gray}{\schulart}}\\[2cm]

    {\Huge\textcolor{fach}{\textbf{Langentwurf}}}\\[0.5cm]
    {\large Tier \plantier{} -- Vollständig \quad|\quad Profil \lerngruppenprofil}\\[2cm]

    {\LARGE\textbf{\stundenthema}}\\[0.5cm]
    {\large im Rahmen der Unterrichtsreihe}\\[0.3cm]
    {\Large\textit{\reihenthema}}\\[2cm]

    \begin{tabular}{rl}
        \textbf{Fach:} & \fach \\[0.3cm]
        \textbf{Klasse:} & \jahrgangsstufe{} (\kursart) \\[0.3cm]
        \textbf{Datum:} & \underline{\hspace{3cm}} \\[0.3cm]
        \textbf{Zeit:} & \underline{\hspace{3cm}} \\[0.3cm]
        \textbf{Raum:} & \underline{\hspace{3cm}} \\[0.3cm]
    \end{tabular}

    \vfill

    {\large Lehrkraft: \lehrkraft}\\[0.3cm]
    {\normalsize\textcolor{gray}{\ausbildungsstand}}\\[1cm]

\end{titlepage}

% ============================================================================
% INHALTSVERZEICHNIS
% ============================================================================
\tableofcontents
\newpage

% ============================================================================
% KONFIGURATIONS-ÜBERSICHT
% ============================================================================
\begin{infobox}[Konfiguration dieses Langentwurfs]
\begin{tabular}{ll}
    \textbf{Tier:} & 1 -- Vollständig (alle Abschnitte) \\[0.3em]
    \textbf{Profil:} & A -- Gesamtschule heterogen \\[0.3em]
    \textbf{Fach:} & Physik \\[0.3em]
    \textbf{Testdifferenzierung:} & Ja (3 Niveaus: */**/***)
\end{tabular}
\end{infobox}

\begin{profilbox}
\profilbeschreibung
\end{profilbox}

% ============================================================================
% 1. BEDINGUNGSANALYSE
% ============================================================================
\section{Bedingungsanalyse}

\subsection{Statistik der Lerngruppe}
\begin{tabular}{ll}
    Kursgröße: & 24 SuS (13 m / 11 w / 0 d) \\
    Jahrgangsstufe: & 8 \\
    Kursart: & Regelunterricht (keine Kursdifferenzierung) \\
    Profil: & A (heterogen) \\
\end{tabular}

\subsection{Zusammensetzung nach Leistungsniveau}
\begin{tabular}{|l|c|l|}
\hline
\textbf{Niveau} & \textbf{Anteil} & \textbf{Charakteristik} \\
\hline
Berufsreife (BR / *) & ca. 5 SuS (20\%) & Qualitative Beschreibungen, einfache Rechnungen \\
\hline
Mittlere Reife (MR / **) & ca. 12 SuS (50\%) & Formelanwendung, Einheiten umrechnen \\
\hline
Gymnasium (GY / ***) & ca. 7 SuS (30\%) & Formelumstellung, Transfer, Fehleranalyse \\
\hline
\end{tabular}

\subsection{Individuelle Besonderheiten (Inklusion)}
\begin{itemize}
    \item \textbf{LRS:} 2 SuS (Nachteilsausgleich: Zeitzugabe 10\%, vergrößerte AB)
    \item \textbf{Förderbedarf Lernen:} 1 SuS (reduzierte Aufgabenmenge, Lückentext statt Freitext)
    \item \textbf{Besonders leistungsstark:} 3 SuS (Zusatzaufgaben, Expertenrolle bei Gruppenarbeit)
    \item \textbf{Verhaltensauffällig:} 1 SuS (klare Strukturen, häufige Aktivierung)
\end{itemize}

\subsection{Arbeitsvoraussetzungen}
\begin{itemize}
    \item Raumsituation: Physikraum mit Experimentierinseln (6 Gruppen à 4 SuS)
    \item Technische Ausstattung: Beamer, Dokumentenkamera, WLAN
    \item Experimentiermaterial: Netzgeräte, Widerstände, Multimeter, Kabel
    \item Digitale Medien: Tablets vorhanden (1 pro Gruppe)
\end{itemize}

\subsection{Fachbezogener Entwicklungsstand}
\begin{itemize}
    \item \textbf{Vorwissen:} Stromkreis, Spannung, Stromstärke (Kl. 7/8), Reihen- und Parallelschaltung
    \item \textbf{Bekannte Formeln:} $I = \frac{Q}{t}$, $U$ als ``Antrieb'' für Elektronen
    \item \textbf{Bekannte Methoden:} Partnerexperimente, Protokollführung, Diagramme erstellen
    \item \textbf{Schwächen:} Einheiten umrechnen (mA $\leftrightarrow$ A), Bruchrechnung
\end{itemize}

\subsection{Erfahrungen mit der Lerngruppe}
\begin{itemize}
    \item Motivation bei Experimenten hoch, bei Rechenaufgaben geringer
    \item Gruppenarbeit funktioniert gut bei klaren Rollenverteilungen
    \item Einige SuS brauchen längere Bearbeitungszeit
    \item Visualisierungen (Schaltbilder, Diagramme) werden gut angenommen
\end{itemize}

% ============================================================================
% 2. SACHANALYSE
% ============================================================================
\section{Sachanalyse}

\subsection{Fachwissenschaftliche Einordnung}
Das \textbf{Ohmsche Gesetz} beschreibt den linearen Zusammenhang zwischen Spannung $U$ und Stromstärke $I$ bei konstantem Widerstand $R$:

\[
U = R \cdot I \qquad \text{bzw.} \qquad R = \frac{U}{I} \qquad \text{bzw.} \qquad I = \frac{U}{R}
\]

\textbf{Einheiten:}
\begin{itemize}
    \item Spannung $U$ in Volt (\si{\volt})
    \item Stromstärke $I$ in Ampere (\si{\ampere})
    \item Widerstand $R$ in Ohm (\si{\ohm}), wobei $\SI{1}{\ohm} = \SI{1}{\volt\per\ampere}$
\end{itemize}

Das Gesetz gilt streng nur für \textbf{ohmsche Widerstände} (metallische Leiter bei konstanter Temperatur). Bauteile wie Glühlampen, Dioden oder Halbleiter zeigen nichtlineares Verhalten.

\subsection{Kontrastive Betrachtung}
\textbf{Typische Fehlvorstellungen:}
\begin{itemize}
    \item ``Der Strom wird im Widerstand verbraucht'' (Stromverbrauchsvorstellung)
    \item ``Höhere Spannung bedeutet immer mehr Strom'' (ohne Berücksichtigung des Widerstands)
    \item ``Der Widerstand ändert sich, wenn ich die Spannung ändere'' (Verwechslung von $R$ und $I$)
    \item Verwechslung von Ursache und Wirkung ($U$ ist Ursache, $I$ ist Wirkung)
\end{itemize}

\textbf{Alltagsvorstellungen nutzen:}
\begin{itemize}
    \item Wassermodell: Druck (Spannung), Durchfluss (Strom), enge Röhre (Widerstand)
    \item Aber: Modell hat Grenzen (Wasser wird ``verbraucht'', Strom nicht)
\end{itemize}

\subsection{Kernkonzepte der Stunde}
\begin{enumerate}
    \item \textbf{Proportionalität:} Bei konstantem $R$ ist $I$ proportional zu $U$
    \item \textbf{Widerstand als Materialeigenschaft:} $R$ ist konstant (bei ohmschen Widerständen)
    \item \textbf{Formel anwenden:} $U = R \cdot I$ in verschiedenen Richtungen umstellen
    \item \textbf{Kopfrechenbar:} Nur ganzzahlige Werte, $R \in \{2, 4, 5, 10, 20, 50, 100\}\,\si{\ohm}$
\end{enumerate}

\subsection{Weiterführende Aspekte}
\begin{itemize}
    \item Stunde 2: Widerstand messen, Kennlinien aufnehmen
    \item Stunde 3: Reihen- und Parallelschaltung von Widerständen
    \item Stunde 4: Anwendungen (Spannungsteiler, Vorwiderstand LED)
\end{itemize}

% ============================================================================
% 3. DIDAKTISCHE ANALYSE
% ============================================================================
\section{Didaktische Analyse}

\subsection{Stellung der Stunde in der Unterrichtseinheit}
Dies ist die \textbf{1. Stunde} der Unterrichtsreihe ``Elektrischer Widerstand'' (4 Stunden gesamt).

\begin{tabular}{|c|l|l|}
\hline
\textbf{Std.} & \textbf{Thema} & \textbf{Schwerpunkt} \\
\hline
\textbf{1} & \textbf{Das Ohmsche Gesetz} & \textbf{Einführung, qualitativ + quantitativ} \\
\hline
2 & Widerstand messen & Kennlinien, ohm'sch vs. nicht-ohm'sch \\
\hline
3 & Widerstände schalten & Reihen- und Parallelschaltung \\
\hline
4 & Anwendungen & Spannungsteiler, Vorwiderstand \\
\hline
\end{tabular}

\subsection{Ziele der Unterrichtseinheit}
Am Ende der Reihe können die SuS...
\begin{itemize}
    \item den elektrischen Widerstand als Materialeigenschaft beschreiben
    \item das Ohmsche Gesetz in verschiedenen Formen anwenden
    \item Gesamt\-wider\-stände in Reihen- und Parallelschaltungen berechnen
    \item technische Anwendungen des Widerstands erklären
\end{itemize}

\subsection{Legitimation}
\textbf{Rahmenplanbezug (M-V, Physik Sek I):}
\begin{itemize}
    \item Kompetenzbereich: Fachwissen (Elektrizitätslehre)
    \item Standard: ``SuS beschreiben den Zusammenhang zwischen Spannung, Stromstärke und Widerstand''
    \item Inhaltsfeld: Elektrischer Stromkreis, elektrischer Widerstand
\end{itemize}

\textbf{Bildungswert:}
\begin{itemize}
    \item \textbf{Gegenwartsbedeutung:} Jedes elektronische Gerät enthält Widerstände
    \item \textbf{Zukunftsbedeutung:} Grundlage für Elektrotechnik-Berufe, Alltagsverständnis
    \item \textbf{Exemplarische Bedeutung:} Proportionale Zusammenhänge, Modellierung durch Formeln
\end{itemize}

\subsection{Didaktische Reduktion}
\textbf{Bewusst ausgeklammert:}
\begin{itemize}
    \item Temperaturabhängigkeit des Widerstands (kommt in Stunde 2)
    \item Spezifischer Widerstand $\rho$ (Oberstufe)
    \item Nichtlineare Widerstände (nur Erwähnung)
    \item Leistung $P = U \cdot I$ (eigene Reihe)
\end{itemize}

\textbf{Begründung:} Fokus auf den zentralen proportionalen Zusammenhang. Überlastung vermeiden.

\subsection{Strukturierung des Unterrichts}
\begin{enumerate}
    \item \textbf{Einstieg (5 min):} Problemfrage ``Warum leuchtet die Lampe unterschiedlich hell?''
    \item \textbf{Erarbeitung (15 min):} Experiment: $U$-$I$-Messung bei konstantem $R$
    \item \textbf{Sicherung I (5 min):} Diagramm erstellen, Proportionalität erkennen
    \item \textbf{Übung (15 min):} Formel anwenden (differenziert)
    \item \textbf{Sicherung II (5 min):} Zusammenfassung, Hefteintrag
\end{enumerate}

% ============================================================================
% 4. STUNDENZIELE
% ============================================================================
\section{Stundenziele}

\subsection{Grobziele}
Die SuS erweitern ihre \textbf{Fachkompetenz im Bereich Elektrizitätslehre}, indem sie den Zusammenhang zwischen Spannung, Stromstärke und Widerstand experimentell erkunden und das Ohmsche Gesetz zur Berechnung anwenden.

\textbf{Kompetenzbereiche:}
\begin{itemize}
    \item \textbf{Fachwissen:} Ohmsches Gesetz verstehen und anwenden
    \item \textbf{Erkenntnisgewinnung:} Experiment durchführen, Daten auswerten
    \item \textbf{Kommunikation:} Ergebnisse in Diagramm darstellen
    \item \textbf{Bewertung:} (in Folgestunden)
\end{itemize}

\subsection{Feinziele (operationalisiert)}
\begin{tabular}{|c|p{8cm}|c|c|}
\hline
\textbf{Nr.} & \textbf{Die SuS können...} & \textbf{AFB} & \textbf{Niveau} \\
\hline
FZ1 & den Zusammenhang ``Je höher $U$, desto höher $I$'' qualitativ \textbf{beschreiben}. & I & alle \\
\hline
FZ2 & Messwerte in ein $U$-$I$-Diagramm \textbf{eintragen} und die Gerade interpretieren. & I/II & alle \\
\hline
FZ3 & die Formel $U = R \cdot I$ \textbf{anwenden}, um $U$, $I$ oder $R$ zu berechnen (bei gegebenen Werten). & II & **/*** \\
\hline
FZ4 & die Formel nach $I$ oder $R$ \textbf{umstellen} und mehrstufige Aufgaben lösen. & II & *** \\
\hline
FZ5 & \textbf{begründen}, warum der Widerstand bei Änderung der Spannung konstant bleibt. & III & *** \\
\hline
\end{tabular}

\vspace{1em}
\textbf{AFB-Verteilung:}
\begin{itemize}
    \item AFB I (Reproduktion): ca. 30\% (FZ1, FZ2 teilweise)
    \item AFB II (Reorganisation/Transfer): ca. 50\% (FZ2 teilweise, FZ3, FZ4)
    \item AFB III (Reflexion/Problemlösung): ca. 20\% (FZ5)
\end{itemize}

% ============================================================================
% 5. METHODISCHE ANALYSE
% ============================================================================
\section{Methodische Analyse}

\subsection{Einstieg}
\begin{itemize}
    \item \textbf{Methode:} Demonstrationsexperiment -- Lampe an Netzgerät, Spannung variieren
    \item \textbf{Impuls:} ``Was passiert mit der Helligkeit? Warum?''
    \item \textbf{Begründung:} Phänomen aktiviert Vorwissen, erzeugt kognitive Neugier
    \item \textbf{Alternative:} Video (falls Experiment nicht möglich)
\end{itemize}

\subsection{Erarbeitung}
\begin{itemize}
    \item \textbf{Methode:} Schülerexperiment in 6 Gruppen (je 4 SuS)
    \item \textbf{Aufgabe:} $I$ messen bei $U = 2, 4, 6, 8, 10\,\si{\volt}$ (Widerstand $R = \SI{20}{\ohm}$)
    \item \textbf{Begründung:} Handlungsorientierung, selbstständige Erkenntnisgewinnung
    \item \textbf{Differenzierung:} Starke SuS helfen schwächeren (Expertenrolle)
\end{itemize}

\subsection{Übungsphasen}
\begin{itemize}
    \item \textbf{Methode:} Differenzierte Arbeitsblätter (3 Niveaus)
    \item \textbf{Progression:} Von Einsetzen (*) über Umstellen (**) zu Transfer (***)
    \item \textbf{Begründung:} Jeder arbeitet auf seinem Niveau, Erfolgserlebnisse für alle
\end{itemize}

\subsection{Sozialformen}
\begin{tabular}{|l|l|l|}
\hline
\textbf{Phase} & \textbf{Sozialform} & \textbf{Begründung} \\
\hline
Einstieg & Plenum & Gemeinsamer Fokus, Aktivierung \\
\hline
Erarbeitung & Gruppenarbeit (4er) & Kooperatives Experimentieren \\
\hline
Sicherung I & Plenum & Gemeinsame Diagrammerstellung \\
\hline
Übung & Einzelarbeit & Individuelle Vertiefung \\
\hline
Sicherung II & Plenum & Zusammenfassung, Klärung \\
\hline
\end{tabular}

\subsection{Medien}
\begin{itemize}
    \item Tafel/Whiteboard: Schaltbild, Diagramm, Formel
    \item Dokumentenkamera: Schülerergebnisse zeigen
    \item Arbeitsblätter: AB-01-ohm.pdf (3 Niveaus)
    \item Experimentiermaterial: Netzgerät, Widerstand (\SI{20}{\ohm}), Multimeter, Kabel
\end{itemize}

\subsection{Differenzierung}
\begin{tabular}{|l|p{3.5cm}|p{3.5cm}|p{3.5cm}|}
\hline
& \textbf{Basis (*)} & \textbf{Standard (**)} & \textbf{Erweitert (***)} \\
\hline
\textbf{Inhalt} & Qualitativ: ``Je mehr $U$, desto mehr $I$'' & Formel einsetzen: $U = R \cdot I$ & Formel umstellen, mehrstufig \\
\hline
\textbf{Aufgaben} & Lückentext, Zuordnung, Einsetzen & Berechnen mit gegebener Formel & Umstellen, Fehlersuche \\
\hline
\textbf{Hilfen} & Formelkarte, Beispielrechnung & Teillösung, Tipps & Nur Impulse \\
\hline
\textbf{Zahlen} & $R \in \{10, 20\}$, $U, I$ ganzzahlig & $R \in \{5, 10, 20, 50\}$ & Alle kopfrechenbaren Werte \\
\hline
\end{tabular}

\subsection{Erwartungshorizonte}
\begin{itemize}
    \item \textbf{Minimum (*):} Qualitative Beschreibung, Diagramm ablesen, einfaches Einsetzen
    \item \textbf{Regelstandard (**):} Formel anwenden, Einheiten korrekt, 3 Berechnungen
    \item \textbf{Optimal (***):} Formel umstellen, Fehler erkennen, Transfer auf neue Situation
\end{itemize}

\subsection{Überprüfung der Teilziele}
\begin{tabular}{|c|l|l|}
\hline
\textbf{FZ} & \textbf{Überprüfung durch} & \textbf{Zeitpunkt} \\
\hline
FZ1 & Mündliche Beiträge im Plenum & Einstieg \\
\hline
FZ2 & Kontrolle der Gruppendiagramme & Sicherung I \\
\hline
FZ3 & Arbeitsblatt-Kontrolle & Übungsphase \\
\hline
FZ4 & Arbeitsblatt (***-Aufgaben) & Übungsphase \\
\hline
FZ5 & Zusatzfrage für Schnelle & Ende \\
\hline
\end{tabular}

% ============================================================================
% 6. VERLAUFSPLANUNG
% ============================================================================
\section{Verlaufsplanung}

\begin{longtable}{|p{1cm}|p{1.5cm}|p{4cm}|p{3cm}|p{2cm}|p{2cm}|}
\hline
\textbf{Zeit} & \textbf{Phase} & \textbf{Inhalt / Lehrerhandeln} & \textbf{Schülerhandeln} & \textbf{Sozial-form} & \textbf{Medien} \\
\hline
\endhead

0--5 & Einstieg & L zeigt Demo: Lampe an Netzgerät, variiert $U$. Impuls: ``Was passiert? Warum?'' & SuS beobachten, äußern Vermutungen & Plenum & Netzgerät, Lampe \\
\hline

5--10 & Über-leitung & L erklärt Experiment: $I$ messen bei verschiedenen $U$. Zeigt Schaltbild. & SuS hören zu, fragen nach & Plenum & Tafel, Schaltbild \\
\hline

10--25 & Erar-beitung & L geht herum, unterstützt. Rollenverteilung in Gruppen. & SuS experimentieren: Messen $I$ bei $U = 2, 4, 6, 8, 10\,\si{V}$ & GA (4er) & Experimentier-material \\
\hline

25--30 & Siche-rung I & L sammelt Werte, erstellt gemeinsam Diagramm. Frage: ``Was fällt auf?'' & SuS nennen Werte, erkennen Proportionalität & Plenum & Tafel, Doku-Kamera \\
\hline

30--40 & Übung & L verteilt differenzierte AB. Geht herum, gibt Hilfen. & SuS bearbeiten AB auf eigenem Niveau & EA & AB-01-ohm (*/**/***) \\
\hline

40--45 & Siche-rung II & L bespricht Musterlösungen. Hefteintrag: Formel + Beispiel. & SuS vergleichen, übernehmen Merksatz & Plenum & Tafel, Heft \\
\hline
\end{longtable}

\textbf{Legende:} EA = Einzelarbeit, GA = Gruppenarbeit

% ============================================================================
% 7. REFLEXION (nach Durchführung)
% ============================================================================
\section{Reflexion (nach der Durchführung)}

\textit{Dieser Abschnitt wird nach der Unterrichtsdurchführung ausgefüllt.}

\subsection{Realisierung der Lernziele}
\begin{tabular}{|c|c|l|}
\hline
\textbf{Feinziel} & \textbf{Erreicht?} & \textbf{Evidenz/Anmerkung} \\
\hline
FZ1 & $\Box$ ja / $\Box$ teilweise / $\Box$ nein & \\
\hline
FZ2 & $\Box$ ja / $\Box$ teilweise / $\Box$ nein & \\
\hline
FZ3 & $\Box$ ja / $\Box$ teilweise / $\Box$ nein & \\
\hline
FZ4 & $\Box$ ja / $\Box$ teilweise / $\Box$ nein & \\
\hline
FZ5 & $\Box$ ja / $\Box$ teilweise / $\Box$ nein & \\
\hline
\end{tabular}

\subsection{Wirksamkeit des Vorgehens}
\begin{itemize}
    \item Was hat gut funktioniert?
    \item Was hat nicht funktioniert?
    \item Unvorhergesehene Ereignisse:
\end{itemize}

\subsection{Konsequenzen für die Weiterführung}
\begin{itemize}
    \item Für die nächste Stunde:
    \item Für ähnliche Stunden:
\end{itemize}

\subsection{Gewonnene Einsichten}
\begin{itemize}
    \item Fachdidaktisch:
    \item Methodisch:
    \item Pädagogisch:
\end{itemize}

% ============================================================================
% 8. ANLAGEN
% ============================================================================
\section{Anlagen}

\begin{enumerate}
    \item Arbeitsblatt (AB-01-ohm.pdf) -- 3 Niveaus
    \item Musterlösung (ML-01-ohm.pdf)
    \item Schaltbild (Tafelanschrieb)
    \item Experimentieranleitung
\end{enumerate}

% ============================================================================
% 9. LÖSUNGEN
% ============================================================================
\section{Lösungen}

\subsection{Lösungen Arbeitsblatt}

\textbf{Niveau * (Basis):}
\begin{enumerate}
    \item Lückentext: ``Je höher die \underline{Spannung}, desto höher die \underline{Stromstärke}.''
    \item $U = R \cdot I = \SI{20}{\ohm} \cdot \SI{0,5}{\ampere} = \SI{10}{\volt}$
    \item Diagramm: Gerade durch Ursprung
\end{enumerate}

\textbf{Niveau ** (Standard):}
\begin{enumerate}
    \item $U = R \cdot I = \SI{50}{\ohm} \cdot \SI{0,2}{\ampere} = \SI{10}{\volt}$
    \item $I = \frac{U}{R} = \frac{\SI{12}{\volt}}{\SI{4}{\ohm}} = \SI{3}{\ampere}$
    \item $R = \frac{U}{I} = \frac{\SI{10}{\volt}}{\SI{0,5}{\ampere}} = \SI{20}{\ohm}$
\end{enumerate}

\textbf{Niveau *** (Erweitert):}
\begin{enumerate}
    \item $R = \frac{U}{I} = \frac{\SI{6}{\volt}}{\SI{0,3}{\ampere}} = \SI{20}{\ohm}$
    \item Fehlersuche: $I = \SI{200}{\milli\ampere} = \SI{0,2}{\ampere}$ (nicht 200!)
    \item Transfer: Bei doppelter Spannung fließt doppelter Strom (Proportionalität)
\end{enumerate}

% ============================================================================
% 10. LITERATUR
% ============================================================================
\section{Literatur}

\subsection{Fachliteratur}
\begin{itemize}
    \item Meschede, D. (2015): Gerthsen Physik. 25. Aufl., Springer.
\end{itemize}

\subsection{Didaktische Literatur}
\begin{itemize}
    \item Mikelskis, H. (Hrsg.) (2006): Physik-Didaktik. Cornelsen.
    \item Kircher, E. et al. (2020): Physikdidaktik. 4. Aufl., Springer.
\end{itemize}

\subsection{Rahmenplan}
\begin{itemize}
    \item Ministerium für Bildung, Wissenschaft und Kultur M-V (2020): Rahmenplan Physik für die Jahrgangsstufen 7--10. Schwerin.
\end{itemize}

% ============================================================================
% 11. DIDAKTIK-SCORE
% ============================================================================
\newpage
\section{Didaktik-Score}

\begin{scorebox}
\textbf{Bewertung der didaktischen Qualität nach 10 Dimensionen}\\
\textbf{Ziel-Score: $\geq$ 9.0 (Premium-Qualität)}

\vspace{1em}
\begin{tabular}{|c|l|c|c|p{4.5cm}|}
\hline
\textbf{\#} & \textbf{Dimension} & \textbf{Gew.} & \textbf{Score} & \textbf{Notiz} \\
\hline
1 & Differenzierung & 15\% & 9/10 & 3 Niveaus konsequent, Inklusion berücksichtigt \\
\hline
2 & Taxonomie (AFB) & 10\% & 9/10 & 30/50/20, Operatoren korrekt \\
\hline
3 & Lernziel-Alignment & 10\% & 10/10 & Exakt am Rahmenplan M-V \\
\hline
4 & Aktivierung & 10\% & 9/10 & Experiment, nicht nur Rezeption \\
\hline
5 & Scaffolding & 10\% & 9/10 & Gestufte Hilfen, Formelkarten \\
\hline
6 & Feedback-Qualität & 10\% & 8/10 & Musterlösungen vorhanden, aber kein digitales Feedback \\
\hline
7 & Sprachliche Klarheit & 10\% & 9/10 & Altersgerecht, Fachbegriffe erklärt \\
\hline
8 & Motivation \& Relevanz & 10\% & 9/10 & Alltagsbezug (Elektronik), Experiment \\
\hline
9 & Inklusion (UDL) & 10\% & 9/10 & LRS berücksichtigt, multiple Repräsentationen \\
\hline
10 & Praktikabilität & 5\% & 10/10 & 45 Min. realistisch, Material vorhanden \\
\hline
\multicolumn{3}{|r|}{\textbf{GESAMT:}} & \textbf{9.1/10} & \\
\hline
\end{tabular}

\vspace{1em}
\textbf{Bewertungsschwellen:}
\begin{itemize}
    \item[\checkmark] \textcolor{successgreen}{$\geq$ 9.0: \textbf{FREIGABE} (Premium-Qualität)}
    \item[$\Box$] \textcolor{warningorange}{8.0--8.9: Iteration empfohlen}
    \item[$\Box$] 6.0--7.9: Überarbeitung erforderlich
    \item[$\Box$] $<$ 6.0: Grundlegende Neukonzeption
\end{itemize}
\end{scorebox}

\subsection{Score-Begründungen}

\textbf{1. Differenzierung (15\%) -- 9/10}
\begin{itemize}
    \item 3 Niveaus mit qualitativ verschiedenen Anforderungen
    \item Inklusion (LRS, Förderbedarf) explizit berücksichtigt
    \item Expertenrolle für leistungsstarke SuS
\end{itemize}

\textbf{2. Taxonomie/AFB (10\%) -- 9/10}
\begin{itemize}
    \item Verteilung 30/50/20 (leicht über AFB I wegen Experiment)
    \item Operatoren: beschreiben, eintragen, anwenden, umstellen, begründen
\end{itemize}

\textbf{3. Lernziel-Alignment (10\%) -- 10/10}
\begin{itemize}
    \item Direkt aus Rahmenplan M-V abgeleitet
    \item Kompetenzbereich Fachwissen + Erkenntnisgewinnung
\end{itemize}

\textbf{4. Aktivierung (10\%) -- 9/10}
\begin{itemize}
    \item Schülerexperiment: Handlungsorientiert
    \item Nicht nur Rezeption, sondern eigene Messwerte
\end{itemize}

\textbf{5. Scaffolding (10\%) -- 9/10}
\begin{itemize}
    \item Gestufte Hilfen: Formelkarte, Teillösung, Impulse
    \item Keine Lösungsverraten
\end{itemize}

\textbf{6. Feedback-Qualität (10\%) -- 8/10}
\begin{itemize}
    \item Musterlösungen vorhanden
    \item Verbesserungspotential: Digitaler Lernpfad mit automatischem Feedback
\end{itemize}

\textbf{7. Sprachliche Klarheit (10\%) -- 9/10}
\begin{itemize}
    \item Klasse 8: Angemessene Satzlänge
    \item Fachbegriffe (Widerstand, Proportionalität) erklärt
\end{itemize}

\textbf{8. Motivation \& Relevanz (10\%) -- 9/10}
\begin{itemize}
    \item Alltagsbezug: Elektronische Geräte
    \item Experiment erzeugt Neugier
\end{itemize}

\textbf{9. Inklusion/UDL (10\%) -- 9/10}
\begin{itemize}
    \item Multiple Repräsentationen: Experiment + Diagramm + Formel
    \item LRS: Zeitzugabe, größere Schrift
    \item Förderbedarf: Reduzierte Aufgabenmenge
\end{itemize}

\textbf{10. Praktikabilität (5\%) -- 10/10}
\begin{itemize}
    \item 45 Minuten realistisch
    \item Experimentiermaterial im Physikraum vorhanden
\end{itemize}

\subsection{Verbesserungsmaßnahmen}
Score $\geq$ 9.0 erreicht -- keine dringenden Maßnahmen.

\textbf{Optionale Verbesserungen:}
\begin{enumerate}
    \item \textbf{Feedback-Qualität:} Digitalen Lernpfad mit automatischem Feedback ergänzen
    \item \textbf{Aktivierung:} Simulation als Ergänzung zum Experiment
\end{enumerate}

% ============================================================================
% ENDE
% ============================================================================
\end{document}
